% ============================================================================
\clearpage
\section*{v4.9.0-r69 Source-Only: The OP5 Research Series of 25 to 27 October}
\addcontentsline{toc}{section}{v4.9.0-r69 Source-Only: The OP5 Research Series of 25 to 27 October}
\label{sec:v490-r69-op5-series}
% ============================================================================

\paragraph{Purpose.}
Between 25 and 27 October 2026 the author worked on OP5 with another language model.  The results were committed to the repository folder
\texttt{analysis/} (six commits, \texttt{3fa5ba9} to \texttt{337d9cb}, 173 files: notes, check scripts, their outputs and a literature catalogue).
The author asked for this work to be integrated.  This layer records what the series establishes, what it leaves open, and three errors it found in
the main body.  The notes themselves stay in \texttt{analysis/}; the layer names the note for each item.  It makes no prediction, writes no status
note and rewrites no earlier layer.

\noindent\textbf{Source-only status.}
This v4.9.0-r69 layer is source-only.  It does not build a PDF, does not assign a new DOI,
and does not authorize a public release.  The v4.9.0 Version DOI
\href{https://doi.org/10.5281/zenodo.23534954}{10.5281/zenodo.23534954} and the Concept DOI
\href{https://doi.org/10.5281/zenodo.19210728}{10.5281/zenodo.19210728} are unchanged by it.

\paragraph{How the series was checked.}
All 23 new check scripts were rerun with SymPy and NumPy; all finish without error, and the regenerated outputs differ from the committed ones only in the
last floating-point digit.  Nine statements were recomputed with newly written code (\texttt{llm/k106\_nachrechnung\_op5.py}, output
\texttt{llm/k106\_out.txt}); all nine hold.  The notes marked ``Independent'' in the series were written by sub-agents of the same model; there is no
external review.

\subsection*{v4.9.0-r69 results}
\label{subsec:v490-r69-results}

\paragraph{Theoremlet v4.9.0-r69.1 (frames, neutrality and the absolute term).}
\begin{enumerate}[leftmargin=2em]
\item[(a)] \emph{A frame change cannot remove a balance defect.}  Under a passive change of frame with frequency factor \(a>0\) and declared amplitude
factor \(b\neq0\), the defect \(\Delta=I-B\) of a balance of the type of Definition~\ref{def:tr-balance} becomes \((b/a)\Delta\); the statement
\(\Delta=0\) is invariant.  Replacing \(\gL\) by \(e^{\eta}\) in frequency, cutoff or both leaves the historical defect between \(-1.0008\) and \(-1.0038\)
(\texttt{OP5\_Frame\_Recheck.md}).
\item[(b)] \emph{Three different neutral objects.}  (i) The midpoint \(\eta_N=(\eta_1+\eta_2)/2\) of two scale frames with offsets \(\pm\delta\); no input frame
need be photon-like.  (ii) On the real state disk the equal mixture of \(\rho\) and its native counterpart \(\sigma_y\rho\,\sigma_y\) (the tachyon map
\(z\mapsto iz\) of v4.9.0-r28) is \(I/2\); with \(J_B=\sigma_x\) the component \(b_x\) survives.  \(I/2\) is a mixture, not a frame change, and carries the label P,
while the old ``photon frame'' is the pure reference state R (v4.9.0-r24).  (iii) Two positive non-parallel null streams select a unique flux-free timelike
observer; there the stress is \(\mathrm{diag}(\rho,\rho,0,0)\), not the vacuum form \(\mathrm{diag}(\rho,-\rho,-\rho,-\rho)\).  The three must not be identified
(\texttt{OP5\_Lambda\_Neutral\_Reassessment.md}, \texttt{OP5\_Neutral\_Frame\_Candidates.md}).
\item[(c)] \emph{\(\bsym^2\) as a relative separation (offered, not adopted).}  With \(\chi=\eta_1-\eta_2\), \(q_{\rm rel}=\tanh^2(\chi/2)\) is invariant under a
common shift and exchange; at \(\pm\eta\) it equals \(\tanh^2\eta=(\cosh2\eta-1)/(\cosh2\eta+1)\), the form of \(\bsym^2\) when \(\gL=\cosh2\eta\).  A potential
\(V=M^4\tanh^2(\chi/2)\) has \(V'/M^4\approx0.384\) at the historical separation, so it gives no static vacuum there.
\item[(d)] \emph{The absolute vacuum term is not selected.}  A shift \(H\mapsto H+C\,I\) of the whole generator changes no probability, gap or beat period
in fixed-background unitary dynamics; stationarity at fixed metric is blind to a constant volume density; in the alternative affine parent action of
\texttt{HC\_Two\_Frame\_Action\_Bridge.md} one local branch solves all field equations for every constant \(C\) (branch-local screening, not selection);
one neutral pair configuration is compatible with different constant vacuum energies (\texttt{OP5\_Absolute\_\allowbreak Selection\_\allowbreak Independent.md},
\texttt{OP5\_Affine\_\allowbreak Vacuum\_\allowbreak Response\_\allowbreak Independent.md}, \texttt{OP5\_Neutral\_\allowbreak Pair\_\allowbreak Dynamics\_\allowbreak Result.md}).
\item[(e)] \emph{Readout protocol and energy.}  A passive change of representation changes no prediction; a different active readout (coupling, order,
time, retained branch) can change states and exchanged energy.  Energy then moves between system, apparatus and bath; none is produced by naming a
frame (\texttt{OP5\_Readout\_Protocol\_Energy\_Result.md}).
\item[(f)] \emph{Doublet selection.}  No nontrivial orthogonal projector on the doublet of v4.9.0-r05 respects both the structural involution and the
corrected M/T exchange; only \(0\) and \(I\) do.  The two ports survive as an exchanged pair (\texttt{OP5\_Symmetry\_Sector\_Selection\_Independent.md}).
\end{enumerate}

\paragraph{Theoremlet v4.9.0-r69.2 (cones, coframes and the cross readout).}
\begin{enumerate}[leftmargin=2em]
\item[(a)] \emph{The probe cone comes from the Clifford structure.}  If the coefficients of the probe equation of \S\ref{sec:field-eq} are local
multiplications, the principal symbol is \(\Gamma^\mu\xi_\mu\) with square \(g^{\mu\nu}\xi_\mu\xi_\nu\,I\); mass, \(M_{16}\), triadic and HC terms do not
change the cone.  A conformal rescaling \(E\mapsto aE\) keeps every cone and changes the volume weight by \(a^4\), so cone data fix neither an absolute
scale nor the weight of a vacuum term (\texttt{OP5\_Null\_Transport\_And\_Signal\_Locality.md}).
\item[(b)] \emph{Two complex probes.}  Two free massless spinor fields give the rank-four coframe \(E=\mathrm{diag}(3,-1,-1,-1)\), \(g=\mathrm{diag}(-9,1,1,1)\);
the readout cone differs from the reference cone, and the induced equation leaves a residual \(8i\,e_1\).  For a constant frame and equal induced
massless (or equal-mass) equations the trace \(e_a^{\ \mu}E^a_{\ \mu}\) would have to be \(0\) and is \(4\) (\texttt{OP5\_Two\_Complex\_Probes.md}).  An exact coupled
exchange solution with \(M=8\sigma_y\) matches only at one point: on the time axis the induced residual is \(544\,i\,t^2u+O(t^4)\) and the readout speed
\(1-80t^2+O(t^4)\) (\texttt{OP5\_Two\_Probe\_Exchange\_Evolution.md}).
\item[(c)] \emph{Two flat metrics.}  Between the reference metric and the readout metric the relative connection is nonzero, both curvatures vanish
(\texttt{OP5\_Relative\_Connection\_And\_Free\_Paths.md}).
\item[(d)] \emph{Connected cross readout.}  If \(E^a=\tfrac12D_\omega j_S^a\) with a parallel selector, then \(T^a=\tfrac12R^a{}_b\,j_S^b\) exactly.  Torsion-free
connection plus an ordinary Einstein vacuum equation \(\mathrm{Ric}=\Lambda_{\rm geom}g\) forces \(\Lambda_{\rm geom}j=0\), so \(\Lambda_{\rm geom}=0\) wherever
\(j\neq0\).  This is an exclusion under these premises, not a cancellation of vacuum energy; with \(a=\tau\) the comoving fluid needs
\(\rho+3p=2\Lambda_{\rm geom}/\kappa\).  A torsionful witness with \(E\) Minkowski exists
(\texttt{OP5\_Cross\_Connection\_Integrability.md}, \texttt{OP5\_Torsionful\_Geometric\_Witness.md}).
\item[(e)] \emph{Ward exclusion.}  For the common-boost triad current and an isotropic time-only boost connection \(h(t)\), current conservation in the
reconstructed geometry requires \(b+3tb'=1\) and \((hb^2)'=0\); the second gives \(b=(1\pm\sqrt{10})/9\), the first \(b=1\).  No differentiable profile exists
(\texttt{OP5\_Isotropic\_Connection\_Ward\_NoGo.md}).
\item[(f)] \emph{Counterflow and chiral tangent.}  A counterflow configuration that meets every current projection still leaves full spinor residuals with norms
between \(3.9\) and \(9.5\) at two sample points.  Adding independent boosts and one chiral tangent \(e^{i\chi\gamma_5/2}\) with signs \((+,+,-)\) fits the
instantaneous derivative at both points (rank \(22\) of \(24\), residual below \(4\cdot10^{-28}\)); this is a two-point fit, not a solution.  The prescribed
state set is not invariant under the comparison dynamics; a rotated probe pair cancels the axial drift but not the radial drift, and
\(\overline{E^T\eta E}\neq\bar E^T\eta\bar E\) in general (\texttt{OP5\_Counterflow\_Full\_Residual\_Result.md}, \texttt{OP5\_Chiral\_Tangent\_Result.md},
\texttt{OP5\_Constraint\_Set\_And\_Ensemble\_Result.md}).
\end{enumerate}

\paragraph{Finding v4.9.0-r69.3 (errors in the main body).}
\begin{enumerate}[leftmargin=2em]
\item[(a)] Remark~\ref{rem:regularization-B}: the expansion is \(R(\tau)=-1/\tau+1-\tfrac{11}{24}\tau+\tfrac18\tau^2+O(\tau^3)\), not
\(-1/\tau+\tfrac12-\tau/12\).  Since \(\int_\varepsilon^T(1+R)\,d\tau\) already contains \(+\log\varepsilon\), the printed counterterm \(+\log\varepsilon\) doubles the
divergence; the counterterm is \(-\log(\varepsilon/\mu)\) with a declared scale \(\mu\).  With \(\mu=1\) and \(T=0.393\) the value is \(1.689\), not \(0.940\), so
the \(0.893\) of Proposition~\ref{prop:theta-required-B} is not supported.  The route was already marked false (v4.8.0-r05); the numbers are new.
\item[(b)] Theorem~\ref{thm:dirac-probe} and Proposition~\ref{prop:three-compatible-readings} still fix \(\bsym\) and the clock scale by \(\gL^*\) of
Theorem~\ref{thm:gue-selfdet}, which is archived.  The citation no longer derives a clock scale.
\item[(c)] The corollary ``Cosmological constant revisited'' of \S\ref{sec:phase-constants} calls the photon-frame term zero and gives the finite ratio
\(\rho_{\rm vac}^{(\rm photon)}/\rho_{\rm vac}^{(\rm matter)}=\gL^5\).  For one and the same quantity both cannot hold; a distinction between net and bare
density would have to be defined.  This adds to the status-note candidates of v4.9.0-r64 and r67.
\end{enumerate}
(a) to (c) are candidates for status notes (rules B and C of v4.9.0-r25); they are not written in this layer.

\paragraph{Remarks v4.9.0-r69.4.}
\begin{enumerate}[leftmargin=2em]
\item The series corrects itself once: \texttt{OP5\_Comparison\_Map\_And\_Frame\_Type.md} called the realification of \(J_B\) missing; it is
Theoremlet v4.9.0-r04.1 (\texttt{OP5\_Realification\_And\_Gram\_Boundary.md}).
\item Every positive result of the series is conditional: a declared comparison model, a chosen witness, or two sample points.  None is derived from
the native \(M_{16}\), \(P_{C_7}\) or a common action.  Every note states that OP5 stays open.
\item The literature routes (\texttt{OP5\_Mathematical\_Literature\_Routes\_2026-10-27.md}) suggest, in this order: polar spinor variables
(Fabbri, \texttt{arXiv:1707.03270}) and Fierz reconstruction (Inglis--Jarvis, \texttt{arXiv:1405.6762}); the constraint algorithm for
\(k\)-presymplectic systems (Gracia, Mart\'in, Rom\'an-Roy, \texttt{arXiv:0903.1791}); formal integrability (Landsberg, \texttt{dg-ga/9708005}); then a
coupled coframe dynamics.  Green and Wald (\texttt{arXiv:1011.4920}) and Kaloper et al.\ (\texttt{arXiv:1505.01492}) bound what averaging and sequestering
can do for \(\Lambda\).
\item The clock window \(E_{\rm clk}^{\rm win}\) of the C5 chain is unchanged by a joint relabelling \(\tau\mapsto a\tau\); it selects a relay window, not a
clock period.  This does not touch \(W_{\rm clk}^{\rm loc}\) under \(\Delta_\gamma<1/5\).
\end{enumerate}

\paragraph{Proof sketch.}
r69.1 and r69.2 by the notes named in each item and their check scripts in \texttt{analysis/}.  Recomputed with new code
(\texttt{llm/k106\_nachrechnung\_op5.py}): r69.1(b)(ii) and (c) including \(V'\), r69.2(b), (d) and (e), and the series and value of r69.3(a).  r69.3 also by reading the
cited passages.
\(\square\)

\subsection*{v4.9.0-r69 status in the three frames}
\label{subsec:v490-r69-frames}

\begin{itemize}[leftmargin=2em]
\item r69.1(a) holds in every frame: the defect transforms by \(b/a\) between M, T and R alike.  It concerns finite invertible frame changes; the
photon P is a degenerate limit and needs its own normalized observable.
\item r69.1(b): (i) needs two real scale frames, so it compares M with R or with \(T|_R\), not with the native T (complex rapidity, v4.9.0-r24); (ii) is built
from M and its T-image and gives the P state; (iii) is a spacetime observer and has no frame label.  Frame-free only in the sense that (ii) is symmetric
in M and T.
\item r69.1(d), (e), (f) and r69.2 are statements in the reference description R or in a declared classical comparison sector; they do not separate M
from T.
\item r69.3 concerns R and P: the old ``photon frame'' of the cited passages is R (v4.9.0-r24).
\end{itemize}
No statement of this layer concerns colour or charge.  The chiral phase \(e^{i\chi\gamma_5/2}\) acts in the Dirac comparison sector; it is neither the inner
unit \(J=L_{e_7}\) nor the \(\mathbb C\)-factor of the Dixon register.

\subsection*{v4.9.0-r69 register entries}
\label{subsec:v490-r69-register}

\begin{center}
\small
\begin{tabular}{@{}>{\raggedright\arraybackslash}p{0.34\textwidth}>{\raggedright\arraybackslash}p{0.58\textwidth}@{}}
\toprule
statement & status after v4.9.0-r69\\
\midrule
OP5 research series (\texttt{analysis/}, 25--27 Oct.) & reproducible; conditional results; OP5\(_{\rm exact}\) open\\
frame change as a repair of the \(\Theta\) balance & excluded for finite invertible changes\\
neutral reference & three distinct objects; none fixes the absolute vacuum term\\
Remark~\ref{rem:regularization-B} & expansion and counterterm wrong; status-note candidate\\
\(\gL^*\) clock scale in Thm.~\ref{thm:dirac-probe}, Prop.~\ref{prop:three-compatible-readings} & rests on an archived theorem; status-note candidate\\
corollary ``Cosmological constant revisited'' & zero and \(\gL^5\) for one quantity; status-note candidate\\
\bottomrule
\end{tabular}
\end{center}

\subsection*{v4.9.0-r69 ledger and guard update}
\label{subsec:v490-r69-ledger-guard}

\[
R_{\rm chk}^{v490\text{-}r69}=(1,1,0)
\]
for (series reproduced and nine statements confirmed; three errors in the main body found; OP5 not closed).

The layer spends no guard:
\[
R_{\rm glob}^{(v490\text{-}r69,+)}=(0,1,1,1,1),
\qquad
\Gamma_{\rm guard}^{v490\text{-}r69}=(1,1,1,1,1),
\qquad
\beta_4=1/256\ \text{unspent}.
\]
It does not imply exact finite \(C5\) zero closure, exact all-readout terminality,
\(C6/T2\), a global photon-clock/RH theorem, or global all-readout closure.
